\section{Summary}

In this paper, we showed that quasi-PDFs may be seen as hybrids of PDFs and the
primordial rest-frame
momentum distributions of partons.
In this context, the parton's $k_3$ momentum
comes from the motion
of the hadron as a whole and from the primordial
rest-frame
momentum distribution. The complicated convolution nature of quasi-PDFs
necessitates using $p_3 \gtrsim $ \mbox{3 GeV} to wipe out
the primordial momentum distribution effects and
get reasonably close to the PDF limit.

As an alternative approach, we propose to use
pseudo-PDFs ${\cal P} (x, z_3^2)  $ that generalize
the light-front PDFs onto spacelike intervals. By a Fourier transform,
they are related to the
Ioffe-time distributions ${\cal M} (\nu, z_3^2) $ given by generic matrix elements
written as functions of $\nu = p_3 z_3$ and $z_3^2$.
The advantageous features of pseudo-PDFs are that they, first, have the
same $-1 \leq x \leq 1$ support as PDFs, and second, their
$z_3^2$-dependence for small $z_3^2$ is governed by a usual evolution equation.

Forming the ratio ${\cal M} (\nu, z_3^2)/{\cal M} (0, z_3^2)$
of Ioffe-time distributions one
divides out the bulk of $z_3^2$ dependence generated by
the primordial rest-frame distribution.
Furthermore, taking this ratio
one can exclude the \mbox{$z_3^2$-dependent} factor coming from the
lattice renormalization of the $\hat E (0,z_3;A)$ link
creating difficulties (see, e.g., \cite{Ishikawa:2016znu}) for lattice calculations of quasi-PDFs.

Testing the efficiency of using pseudo-PDFs for lattice extractions
of PDFs is a challenge for future studies.

In fact, while this paper was in the review
process, an actual lattice calculation \cite{Orginos:2017kos} based on the ideas of
the present paper was performed.
It has clearly demonstrated the presence of a linear component in the $z_3$-dependence
of the rest-frame function ${\cal M} (0,z_3^2)$, that may be attributed
to the $Z(z_3^2/a^2) \sim e^{-c|z_3|/a}$ behavior
generated by the gauge link. It was also observed that
the ratio ${\cal M} (Pz_3 , z_3^2)/{\cal M} (0, z_3^2)$
has a Gaussian-type behavior with respect to $z_3$, which indicates that
the $Z(z_3^2/a^2)$ factors entering into the numerator and denominator
of the ${\mathfrak M} (Pz_3, z_3^2)$ ratio have been canceled, as we expected.

Furthermore, it was found that when plotted as a function of $\nu$ and $z_3$,
the data for the reduced distribution ${\mathfrak M} (\nu, z_3^2)$ have a very mild dependence
on $z_3^2$. This observation indicates that
the soft part of the \mbox{$z_3^2$-dependence} of ${\cal M} (\nu, z_3^2)$
has been canceled by the rest-frame density ${\cal M} (0, z_3^2)$.
This phenomenon corresponds to factorization of the $x$- and $k_\perp$-dependence
for the soft part of the TMD ${\cal F} (x, k_\perp^2)$.

It was also demonstrated that the residual \mbox{$z_3$-dependence}
for small $z_3 \leq 4 \,a$,
may be explained by perturbative evolution, with the $\alpha_s$ value
corresponding to $\alpha_s/\pi =0.1$.
